\chapter*{Modalidad de Grado}

Seleccionar una
\\[1cm]
-	Desarrollo de un proyecto investigativo disciplinar o multidisciplinar
Esta modalidad de grado se refiere al desarrollo de proyectos de investigación que respondan a problemáticas concretas en diferentes contextos socioeconómicos y culturales. Se fundamenta en la aplicación de los avances de la ingeniería en diferentes áreas de énfasis para abordar problemáticas sociales de los contextos local, regional, nacional o internacional.
Cuando el trabajo de grado sea en grupo, se permitirá un máximo de tres estudiantes. La dirección del programa designará un tutor que acompañe el proceso desde su formulación hasta la sustentación del informe final.
\\[1cm]
Producto escrito
Esta modalidad de grado incluye la entrega de un informe escrito y su correspondiente sustentación. La valoración del informe debe tener en cuenta la coherencia y el rigor metodológico; la fundamentación teórica; la presentación y análisis de los resultados; las conclusiones y recomendaciones, y la correspondencia entre las necesidades del entorno, la naturaleza del proyecto y los resultados obtenidos.
\\[1cm]
-	Producción intelectual relevante
La relevancia de esta modalidad de grado depende de la calidad de la producción académica. En tal sentido, la naturaleza del producto debe cumplir con claros estándares de calidad, innovación, sistematicidad y correspondencia respecto a los propósitos formativos del programa al que se vincula el estudiante. De esa forma, se asume como producción intelectual relevante la publicación de un libro, capítulo de libro o artículo publicado en una revista, institucional o externa, especializada en los
campos de conocimiento inherentes al programa del estudiante
Estas publicaciones deben contar con los respectivos registros de isbn o issn, dependiendo de si es un libro o parte de una publicación periódica, bien sea en formato digital o impreso. Dentro de esta modalidad también se contemplan las revistas indexadas alojadas en plataformas científicas avaladas a escala internacional.

Asimismo, se contempla como producción intelectual relevante el diseño de un software relacionado con el campo de formación; el diseño o la modelación original de algún tipo de sistema o subsistema ligado con el programa de estudio y validado por un experto; la construcción y validación de nuevos protocolos, y la formulación o realización de proyectos de emprendimiento. La dirección del programa de Maestría en Ingeniería designará un tutor que acompañe el proceso, con el objetivo de garantizar
el desarrollo de las competencias investigativas requeridas y la calidad del producto final.
\\[1cm]
Producto escrito
Asumiendo que la producción intelectual generada por el estudiante no es un simple requisito, en el caso de la publicación de un libro, capítulo de libro o artículo, no será necesario un informe escrito adicional. Para los otros tipos de producto contemplados en esta modalidad, se deberá entregar un informe que contenga los propósitos, procedimientos y resultados que condujeron a la elaboración del producto entregado. Nota: debe tener la afiliación de la Universidad de La Salle.
\\[1cm]
-	Elaboración de una propuesta de política pública en el campo del conocimiento
Se refiere a la posibilidad de construir una propuesta de política pública o de participar en la construcción de esta a partir de una disciplina o campo interdisciplinar inherente al programa de formación del estudiante. En todo caso, la propuesta o generación de política pública debe plantear alternativas en uno o varios ámbitos del desarrollo local, regional o nacional.
Esta política pública debe orientarse por principios de desarrollo social asociados al Estado, no a intereses particulares. Debe explicitar los intereses públicos a los que sirve y precisar por qué son una alternativa al modelo o a los modelos existentes. La propuesta debe ser valorada por un experto en política pública que conozca el campo. Una vez avalada por este, el estudiante o grupo de estudiantes la presentará de manera formal ante las instancias gubernamentales pertinentes. La dirección del programa
designará un tutor que acompañe el pro- ceso, para garantizar el desarrollo de las competencias investigativas requeridas y la calidad del producto final.
\\[1cm]
Producto escrito
Dadas las características de esta modalidad, según el caso, el producto escrito consistirá en la propuesta en sí misma, acompañada por su soporte de radicación ante la entidad correspondiente, o en el documento de la política pública con los soportes que demuestren su ejecución. En cuanto a su estructura, contará con cada uno de los elementos constitutivos de una política pública según el campo correspondiente. Un experto en política pública deberá valorar este producto una vez el tutor de la modalidad emita el aval correspondiente.

